\documentclass[11pt]{article}
\usepackage{setspace}
\usepackage{geometry}
\usepackage{parskip}
\geometry{margin=1in}
\setstretch{1.12}

\begin{document}

\begin{flushleft}
Zhongli Wang\\
Bielefeld, Germany\\
Phone: +49 160 9541 6015\\
Email: nkwangzhongli@gmail.com\\
Website: https://sites.google.com/view/zhongli-wang
\end{flushleft}

\vspace{3mm}

\begin{flushleft}
Recruitment Committee\\
The Wing Lung Bank Building for Business Studies\\
Shaw Campus\\
Hong Kong Baptist University\\
34 Renfrew Road\\
Kowloon Tong, Kowloon\\
Hong Kong
\end{flushleft}

\vspace{0.3cm}

\begin{flushright}
\today
\end{flushright}
\vspace{0.3cm}

\vspace{2mm}

\begin{flushleft}
\textbf{Application for Assistant Professor / Research Assistant Professor in Economics}
\end{flushleft}

\vspace{2mm}

Dear Members of the Search Committee,

I am writing to apply for the position of Assistant Professor / Research Assistant Professor in Economics at the School of Business, Hong Kong Baptist University (HKBU). I am completing my Ph.D.\ in Economics at Bielefeld University and Università Cattolica del Sacro Cuore through the Marie Skłodowska--Curie EPOC programme. My research combines industrial organization, computational economics, and applied econometrics to study how firms, consumers, and investors adapt in data-rich and rapidly changing environments. I am particularly attracted to HKBU’s emphasis on research rigor and its strengths at the intersection of economics and finance, which provide an excellent setting for my work on digital markets, competition, and market integrity.

My dissertation, \textit{Strategic Behavior and Economic Dynamics in ICT Markets: An Agent-based Framework}, analyzes how market outcomes emerge when heterogeneous firms compete under network effects, congestion, and boundedly rational consumer behavior. I develop an agent-based model in which firms adjust pricing, capacity investment, and innovation-related decisions over time while consumers learn from experience and social interactions. The framework produces empirically relevant patterns such as lock-in, market tipping, quality degradation under congestion, and state-dependent strategic conduct. To connect the model to data and to strengthen external validity, I use simulation-based econometrics to estimate key behavioral and technological parameters and to evaluate counterfactual policy scenarios. The broader goal is to provide an empirically grounded and transparent modeling approach for analyzing competition and welfare in technology-intensive markets where institutional and infrastructural frictions matter.

In parallel, I have pursued complementary research streams on financial market behavior, transparency, and AI-enabled research infrastructure. In \textit{Polytope Fraud Theory}, I use supervised machine learning to detect structured irregularities in firm financial statements, with a focus on firms targeted by activist short sellers in U.S.\ markets. The project links behavioral and institutional mechanisms to observable anomalies and proposes a unified perspective on investor protection and market integrity. 
% This work complements my market-design interests by clarifying how information environments shape incentives, monitoring, and strategic manipulation, and it provides tools that can be adapted to auditing, regulation, and risk screening in data-rich settings.
A second strand extends the same computational orientation to the research process itself. In my working paper \textit{Agentic Workflows for Economic Research: Design and Implementation}, I develop a framework for integrating large-language-model (LLM) agents into core stages of economic research---from structured literature mapping and data collection to code generation, robustness checks, and simulation management---while maintaining reproducibility, version control, and transparent documentation. 
% The paper formalizes how specialized agents can be coordinated as a pipeline under human oversight, with clear interfaces for logging decisions and tracking provenance.
% , and reducing common failure modes (e.g., hidden assumptions, non-replicable data steps). 
% This methodological contribution supports my broader agenda by enabling more scalable and auditable empirical and computational workflows, particularly for projects that combine high-dimensional data, text-as-data, and simulation-based counterfactual evaluation.
% Together with my dissertation, these projects form a coherent agenda on computational approaches to applied microeconomics and finance, centered on heterogeneity, dynamic adjustment, and policy-relevant counterfactual analysis.

Methodologically, my comparative advantage lies in building rigorous economic models that remain operational under realistic complexity. I combine structural thinking with scalable computation: agent-based simulation for market dynamics, estimation and validation using simulation-based methods, and machine learning tools for prediction, classification, and efficient approximation (e.g., surrogate modeling). I work primarily in Python and R and have experience with high-performance computing for large-scale simulations and Monte Carlo evaluation. I place strong weight on reproducibility, documentation, and transparent pipelines, which is particularly valuable for research that integrates computation, microdata, and policy evaluation.

Looking ahead, I plan to develop this agenda along three lines. First, I will expand my work on digital competition and market design by studying how platform rules, data governance, and network structure jointly shape market power and consumer welfare. Second, I will extend my research on market transparency by building broader empirical designs around how institutions and information intermediaries affect fraud detection, disclosure incentives, and market efficiency. Third, I aim to connect computational models more tightly to empirical and experimental evidence by using microdata and, when suitable, controlled behavioral measurements to discipline heterogeneity in beliefs, learning, and decision rules. These directions naturally connect to HKBU’s strengths in applied microeconomics and finance, and to its interest in understanding the economic and regulatory implications of digital transformation.

HKBU would be an exciting environment for me not only because of its research profile, but also because of its international orientation and its commitment to high-quality teaching. I would welcome the opportunity to contribute to undergraduate and postgraduate teaching in areas such as industrial organization, digital economics, applied econometrics, computational methods for economists, and topics at the interface of economics and data science. I also value collaborative research cultures and would look forward to contributing to seminars, mentoring, and joint projects within the Department and across HKBU.

Thank you for considering my application. I would be honored to join HKBU and contribute to its research and teaching mission, and I would welcome the opportunity to discuss how my work fits the Department’s priorities.

\vspace{4mm}

Sincerely,\\[10mm]
Zhongli Wang

\end{document}
